\chapter{冪級数関数}
    \section{2次以上の項が\texorpdfstring{$o(x)\;(x\to 0)$}{o(x) (x→0)}であること}
        \begin{shadebox}
            冪級数関数
            $f(x) = \sum_{i=0}^\infty a_i x^i$の収束半径が$r>0$であるとする。
            $x$の2次以上の項は$o(x)\;(x\to 0)$である。
            よって$a_1\neq 0$ならば，十分小さい$x \neq 0$に対して$f(x)$の2次以上の項は1次までの項に対して真に小さい。
        \end{shadebox}
        \begin{proof}
            \quad\par
            $|x|<r$とする。このとき$|x|<u<r$なる$u$が存在する。
            収束半径の定義より$f(u)$は絶対収束するから$\lim_{i\to\infty} |a_i u^i|$なので$\forall i,\; |a_i u^i| < M$なる$M$が存在する。
            $f(x)$の2次以上の項は
            \[ \left|\sum_{i=2}^\infty a_i x^i\right| \leq \sum_{i=2}^\infty |a_i| |x^i| < \sum_{i=2}^\infty \frac{M}{u^i} |x^i| = M\frac{|x|^2/u^2}{1-|x|/u} \]
            であるから$x\neq 0$ならば
            \[ \left|\sum_{i=2}^\infty a_i x^i\right|/|x| < \frac{M}{u}\frac{|x|/u}{1-|x|/u} \to 0\; (\text{as } x\to 0) \]
        \end{proof}